\scanpage{14}

et dans le cas actuel $\bigwedge^2 E\overset{\unsure{\mathrm{déf}|\mathrm{dfn}|\mathrm{df}}}=L$ est inversible, \deleted{et de plus} donc

\[0\mathrel{-}L\longrightarrow E\otimes E\mathrel{-}\mathrm{Sym}^2(E)\mathrel{-}0,\]

d'où en multipliant par $L^\vee$

\[0\longrightarrow\underline{\mathcal O}\longrightarrow\underbrace{(E\otimes E)\otimes L^\vee}_{\rotatebox{90}{$\simeq$}\ E^\vee\otimes E\simeq\underline{\mathrm{End}}(E)}\mathrel{-}\mathrm{Sym}^2(E)\otimes L^\vee\mathrel{-}0\]

Je dis que par l'iso. can. vertical, le sous-Module $L\otimes L^\vee\simeq\underline{\mathcal O}$ est transformé en le sous-Module $\underline{\mathcal O}.1_E$ de $\underline{\mathrm{End}}(E)$, de façon plus précise qu'on a \inserted{un diagramme} commutatif \deleted{de \unsure{du|des}} d'iso. canonique

\diagram{\begin{tikzcd}[ampersand replacement=\&,column sep=small,row sep=large]
0\arrow[r,dash]\&\underset{\underline{\mathcal O}}{\underbrace{\overbrace{L}^{\det E}\otimes L^\vee}}\arrow[r]\arrow[d,"\sim"'," {\mathrm{id}_{\underline{\mathcal O}}}"]\&(E\otimes E)\otimes L^\vee\arrow[r,dash]\arrow[d,"\sim"']\&\overbrace{\mathrm{Sym}^2(E)}^{\mathrm{Quad}(E^\vee)\simeq\mathrm{Quad}(E)\otimes L^{\otimes2}}\otimes L^\vee\arrow[r]\arrow[d,"\sim"']\&0\\
0\arrow[r,dash]\&\underline{\mathcal O}\arrow[r,hook]\&\underset{\substack{\simeq\\\underbrace{E^\vee}_{E\otimes L^\vee}\otimes E}}{\underline{\mathrm{End}}(E)}\arrow[r,dash]\&\underline{\mathrm{End}}(E)/\underline{\mathcal O}\arrow[r,dash]\&0,
\end{tikzcd}}

d'où je retiens pour l'instant surtout

\diagram{\begin{tikzcd}[ampersand replacement=\&]
\mathrm{Sym}^2(E)\otimes L^\vee\simeq\mathrm{Quad}(E)\otimes L\quad(\simeq\mathrm{Quad}(E,L))\arrow[d,"\sim"']\\
\underline{\mathrm{End}}(E)/\mathcal O
\end{tikzcd}}
% archive sheet number: 13
